\subsection{Opened sources for the configuration arguments}

The following records specify the statements actually used. A reference
to a source theorem is distinguished from the additional deduction
made here. All the attached sources in this table were opened in this
writing. The bibliography marks the checked source versions as verified;
the unresolved journal-number cross-reference is marked separately.

\paragraph{SGA 1, Expos\'e III, Theorem 7.3 and Corollary 7.4.}
Theorem 7.3 assumes a complete local ring $A$ with residue field $k$
and a smooth projective $k$-scheme $X_0$ satisfying
\[
 H^2(X_0,T_{X_0/k})=H^2(X_0,\mathcal O_{X_0})=0.
\]
It gives a smooth projective $A$-scheme reducing to $X_0$; the
following paragraph also treats a prescribed smooth truncated lift.
Corollary 7.4 states the smooth proper curve case. These are at
original margin pages 85--86 in the opened SMF source. They support
the curve lift in A.CF3, since a curve has no coherent cohomology in
degree two. The marked point is an additional elementary step:
formal smoothness lifts it through the nilpotent thickenings, and
properness and completeness give the compatible $A$-section. Neither
statement is used in A.DC1--A.DC4.

\paragraph{SGA 1, Expos\'e XIII, Corollary 2.9.}
The source assumes a smooth proper morphism of finite presentation
$X\to S$, its open complement $U$ of a relative normal-crossings
divisor, and the primes invertible at every residue characteristic of
$S$. A finite locally constant constructible sheaf of groups with
those primes, tame relative to the pair, has locally constant
constructible first direct image and is cohomologically proper through
degree one. This is a statement about finite torsors and their base
change. In the strictly henselian lifting situation of A.CF3, apply
it to every constant finite prime-to-$p$ group; the resulting torsor
comparison gives the prime-to-$p$ fundamental-group isomorphism.
This conclusion is a deduction from the quoted cohomological statement,
not a configuration exactness assertion in the source.

\paragraph{SGA 1, Expos\'e XIII, Section 2.10 and Lemma 2.11.}
Section 2.10 constructs specialization by strict localization of the
base and a chosen \'etale path. Changing that path changes the
homomorphism by an inner automorphism. Under the smooth proper-pair
hypotheses of Corollary 2.9 the prime-to-residue-characteristic
comparison is an isomorphism. Lemma 2.11 assumes a proper morphism of
finite presentation with relative normal-crossings boundary, a
specialization of geometric base points, and compatible boundary
points under generization; it says that the image of the first tame
inertia subgroup is an inertia subgroup at the second boundary point.
The opened text is at original margin pages 390--392. In A.CF3 the
point section and diagonal are smooth relative divisors, and the fibre
is transverse with multiplicity one, so the fundamental groups, maps,
and specified inertia are all transported. These specialization
statements are excluded from the direct proof A.DC1--A.DC4.

\paragraph{SGA 1, Expos\'e XIII, Proposition 4.3 and Examples 4.4.}
The notation in Section 4.0 is essential. If
$K=\ker(\pi_1(X)\to\pi_1(S))$, the middle group $\pi_1'(X)$ is
$\pi_1(X)/\ker(K\to K^L)$; it retains the full base group.
Proposition 4.3 assumes connected $S$, a morphism $f:X\to S$ that
is $0$-acyclic and locally $1$-aspherical for $L$, a section of $f$,
cohomological properness in degrees at most one for finite constant
$L$-groups, constructibility of the direct-image torsor stacks, and
the indicated constructibility of direct images of finite covers.
It concludes
\[
 1\to\pi_1(X_{\bar s})^L\to\pi_1'(X)
 \to\pi_1(S)\to1.
\]
Examples 4.4 verify these hypotheses for the proper smooth case and
for a smooth proper compactification with relative normal-crossings
boundary and a section, with primes invertible on the base. The
open-pair case explicitly cites XIII 2.9. Thus the source does not
justify applying $L$ to the whole total and base groups, and does
not allow $p$ in $L$ in characteristic $p$. These points are respected
by the alternative Stix argument. They are also why the direct
central-induction proof was supplied for the literal dependency
requirement in G1(ii).

\paragraph{SGA 1, Expos\'e XIII, Proposition 5.2 and Corollary 5.3.}
Proposition 5.2 assumes a regular local scheme $X$, a normal-crossings
divisor with equations $f_i$ belonging to a regular parameter system,
and a finite \'etale cover of its complement tame along that divisor.
Let $n_i$ be the least common multiple of the inertia orders appearing
above the generic point of its $i$-th component. Then the $n_i$ are
invertible, and after adjoining $T_i$ with $T_i^{n_i}=f_i$ the cover
extends uniquely, up to unique isomorphism, to a finite \'etale cover
of the entire chart. The proof uses normality for uniqueness and
purity for extension. Corollary 5.3 computes the strictly local tame
group as a product of prime-to-residue-characteristic procyclic
inertias. These statements occur at original margin pages 428--431.

No root stack appears in this source. The root-stack formulation in
A.DC4 is proved from these chart statements: extend the cover on
each chart, extend the descent isomorphisms by uniqueness, and descend
the resulting equivariant finite \'etale covers. The boundary is the
root gerbe of the normal line. Its stabilizer is not assumed to inject
into its fundamental group; injectivity of the surviving infinite
$r$-inertia is proved separately by A.DC3 and the configuration fibre
injection. The same local chart, with one relative parameter, proves
smoothness over the second curve in A.DC2.

\paragraph{SGA 1, Expos\'e X, Proposition 1.2.}
For a proper separable morphism to a locally noetherian base, its Stein
factor is finite \'etale. Here ``separable'' means flat with
geometrically reduced fibres (Definition 1.1). This is the statement
at original margin page 262, also identified in its editorial note
with EGA III 7.8.10(i). It applies to the normalization in A.DC2,
which the tame local charts show to be smooth proper over the base.
The result, not a nonexistent ``Corollary 1.3'', supplies the finite
\'etale family of connected components used in that proof.

\paragraph{Stix, the opened Proposition 2.7.}
In the attached arXiv version and the author-hosted text dated
18 April 2005, the hypotheses are a connected locally noetherian
scheme $S$, a smooth proper genus-$g$ curve with $n$ disjoint ordered
sections, $2-2g-n<0$, and a set $L$ of primes invertible on $S$.
For the complement $C^0$ of the sections, the source defines
$\pi_1^{(L)}(C^0)$ by completing only the kernel of
$\pi_1(C^0)\to\pi_1(S)$. Proposition 2.7 gives the exact sequence
\[
 1\to\pi_1(C^0_{\bar s})^L\to\pi_1^{(L)}(C^0)
 \to\pi_1(S)\to1
\]
and identifies it as the pullback of
$1\to\operatorname{Inn}(P)\to\operatorname{Aut}(P)
\to\operatorname{Out}(P)\to1$ through the base monodromy.
Its conjugation map restricts to the inner action of the fibre.
Taking $L=\{r\}$ and the constant compactified deck-graph family
meets every hypothesis, even if the proper deck group has a $p$-part.
The resulting monodromy on mod-$r$ fibre homology is an $r$-group of
shears, so the finite automorphism argument in A.CF1 supplies the
required total pro-$r$ detection. This validates the source input in
the earlier A.CF12 proof. The omitted tame middle-exactness step is
supplied by A.DC2.

The arXiv source and author text both number this statement 2.7.
The published journal has not yet been opened directly: OUP requests
returned HTTP 403. Mochizuki--Tamagawa Remark 2.2.1 cites a related
Stix assertion as Proposition 2.3, but we do not treat that
cross-reference as verification of a journal-number match.
The source version and theorem content used here are verified; that
number correspondence is explicitly unverified. This limitation does
not affect A.DC1--A.DC4, which do not use Stix's theorem.

\paragraph{Bellingeri, Theorem 1.2 and Section 2.1.}
Theorem 1.2 is for the braid group $B(n,F)$ of a closed orientable
surface of genus $g\geq1$. It uses generators
$\sigma_1,\ldots,\sigma_{n-1},a_1,\ldots,a_g,b_1,\ldots,b_g$,
commutator $[a,b]=aba^{-1}b^{-1}$, the Artin braid relations,
relations (R1)--(R4), and the surface relation (TR)
\[
 \prod_{j=1}^g[a_j,b_j^{-1}]
 =\sigma_1\cdots\sigma_{n-2}\sigma_{n-1}^2
   \sigma_{n-2}\cdots\sigma_1.
\]
For $n=2$, (R1) and the ordinary braid relations are empty. There are
$2g$ instances of (R2), $4\binom g2$ of (R3), $g$ of (R4), and
one (TR); for genus two these are eleven relations. A.CF5--A.CF6
use exactly this convention. Section 2.1 states the Fadell--Neuwirth
fibration and its exact homotopy sequence, and observes that the
higher homotopy groups vanish for an orientable surface other than
the sphere. For our $g\geq2$, this gives the discrete short exact
sequence. Interchanging the two coordinates changes the convention
of forgetting the last strand into the second projection used here.
The attached arXiv source calls the closed-surface theorem 1.2.
The journal metadata (\emph{Journal of Algebra} 274(2), 543--563,
DOI 10.1016/j.jalgebra.2003.12.009) was checked on the publisher's
primary page; no unobserved journal re-numbering is asserted.

\paragraph{Mochizuki--Tamagawa, Definition 1.1 and Proposition 2.2.}
Definition 1.1 makes a PT-formation either a primary full formation
(a singleton prime) or the total formation of all finite groups.
Proposition 2.2 assumes connected $S$, a PT-formation $\mathcal C$
invertible on $S$, and, for every $\ell\in\Sigma_{\mathcal C}$,
$\mathcal C$-group images of both the mod-$\ell$ cyclotomic character
and the mod-$\ell$ homology monodromy of the compactified proper
curve. It gives the length-one configuration fibre exact sequence
in (i), finite generation and slimness in (ii), the base exact
sequence in (iii), algebraically closed field-extension invariance
in (iv), and complete-DVR specialization in (v).
For A.CF9 with the total formation over $\mathbb C$, invertibility
holds, the base is a geometric point, and the monodromy conditions
are automatic. Thus (i) supplies precisely the full profinite
exactness used there. The proposition does not assert exactness for
any proper prime set with at least two primes. Its restriction is
therefore retained; the converse classification is proved by
A.CF7--A.CF8, not attributed to the source. The attached author PDF
was opened through its full proof, printed pages 11--14 (journal
pages 85--88 for the edition beginning at page 75).

\paragraph{The additional direct-proof cohomology source.}
We opened and downloaded J.~S.~Milne's \emph{Lectures on \'Etale
Cohomology}, version 2.21, dated 22 March 2013. The relevant numbered
results are Theorem 12.7 (Leray), Propositions 14.1--14.2 (prime-to-
characteristic Picard and curve cohomology), Theorem 14.9 and Remark
14.10 (Cartan--Leray and pro-covers), Theorem 16.7 (purity for a smooth
pair), and Theorem 17.10 (proper base change). Their hypotheses hold
for our smooth varieties over an algebraically closed field with
$r$ invertible. The open-family calculation and nonabelian induction
are proved here from these inputs. The Picard argument uses the
algebraic Jacobian proof, not the illustrative complex-analytic
discussion preceding it in the notes. The downloaded PDF has SHA-256
\begin{center}\footnotesize\ttfamily
ac4f122f371d38a44c58c296b7dbf88081d89d2de2334070bff3606771c01077.
\end{center}
